\documentclass[12pt,a4paper]{article}

% -------------------------------------------------
% Pakete
% -------------------------------------------------
\usepackage[utf8]{inputenc}
\usepackage[T1]{fontenc}
\usepackage[ngerman]{babel}
\usepackage{lmodern}
\usepackage{microtype}
\usepackage{amsmath,amssymb,amsfonts,amsthm,mathtools}
\usepackage{geometry}
\usepackage{graphicx}
\usepackage{booktabs}
\usepackage{hyperref}
\usepackage{xcolor}
\usepackage{enumitem}
\usepackage{csquotes}

\newcommand{\Spec}{\operatorname{Spec}}
\newcommand{\GL}{\operatorname{GL}}
\newcommand{\PGL}{\operatorname{PGL}}
\newcommand{\tr}{\operatorname{tr}}
\newcommand{\trd}{\operatorname{trd}}
\newcommand{\vol}{\operatorname{vol}}
\newcommand{\diag}{\operatorname{diag}}

\geometry{margin=2.7cm}
\hypersetup{
	colorlinks=true,
	linkcolor=blue!50!black,
	citecolor=blue!50!black,
	urlcolor=blue!50!black,
	pdftitle={Quaternionische Quotientengraphen, Hecke-Faltungsoperatoren und spektrale Vorstufen lokaler Langlands-Geometrie},
	pdfauthor={Thomas Hoffbauer}
}

% -------------------------------------------------
% Robuste Grafikeinbindung
% -------------------------------------------------
\DeclareGraphicsExtensions{.pdf,.png,.jpg}
\graphicspath{{./}{fig/}{figs/}{images/}{img/}}
\newcommand{\includegraphicssafe}[2][]{%
	\IfFileExists{#2}{\includegraphics[#1]{#2}}{\fbox{\texttt{Grafik nicht gefunden: #2}}}
}

% -------------------------------------------------
% Theorem-Umgebungen
% -------------------------------------------------
\theoremstyle{plain}
\newtheorem{satz}{Satz}[section]
\newtheorem{lemma}[satz]{Lemma}
\newtheorem{proposition}[satz]{Proposition}
\newtheorem{korollar}[satz]{Korollar}

\theoremstyle{definition}
\newtheorem{definition}[satz]{Definition}
\newtheorem{hypothese}[satz]{Hypothese}

\theoremstyle{remark}
\newtheorem{bemerkung}[satz]{Bemerkung}
\newtheorem{beispiel}[satz]{Beispiel}

% -------------------------------------------------
% Makros
% -------------------------------------------------
\newcommand{\Q}{\mathbb{Q}}
\newcommand{\R}{\mathbb{R}}
\newcommand{\Z}{\mathbb{Z}}
\newcommand{\Fp}{\mathbb{F}_p}
\newcommand{\C}{\mathbb{C}}
\newcommand{\A}{\mathbb{A}}
\newcommand{\OH}{\mathcal{O}_H}
\newcommand{\HN}{\mathbb{H}}
\newcommand{\Nrd}{\operatorname{Nrd}}
\newcommand{\id}{\mathrm{id}}
\newcommand{\elltwo}{\ell^2}
\newcommand{\cH}{\mathcal{H}}
\newcommand{\cT}{\mathcal{T}}
\newcommand{\cU}{\mathcal{U}}

% -------------------------------------------------
% Titel
% -------------------------------------------------
\title{\textbf{Quaternionische Quotientengraphen, Hecke-Faltungsoperatoren\\
		und spektrale Vorstufen lokaler Langlands-Geometrie}\\[0.5em]
	\large Eine mathematisch vertiefte Spektrum-nahe Fassung}
\author{Thomas Hoffbauer}
\date{März 2026}

\begin{document}
	\maketitle
	
	\begin{abstract}
		Dieser Text untersucht eine elementare, aber arithmetisch motivierte Modellwelt:
		endliche Quotienten der Hurwitz-Ordnung in der Hamiltonschen Quaternionenalgebra
		und die daraus hervorgehenden regulären Graphen sowie Faltungsoperatoren.
		Im Ausgangspunkt ist die Konstruktion diskret und vollständig endlich.
		Sobald die Erzeugermengen jedoch nicht beliebig, sondern über reduzierte Normen
		vorgegeben werden, tritt eine formale Nähe zu Hecke-Operatoren hervor.
		
		Die zentrale These des Artikels lautet nicht, dass hier bereits Langlands-Geometrie
		oder gar ein Spezialfall der Arbeiten von Fargues--Scholze vorliegt.
		Die präzisere Aussage ist bescheidener und mathematisch robuster:
		Das Modell realisiert eine endliche Operatorwelt, in der lokale arithmetische
		Übergänge globale Spektren erzeugen. Diese Architektur ist im schwachen,
		methodischen Sinn verwandt mit einem Grundmotiv moderner arithmetischer Geometrie:
		\[
		\text{lokale Daten}
		\longrightarrow
		\text{Operator}
		\longrightarrow
		\text{Spektrum}
		\longrightarrow
		\text{globale Struktur}.
		\]
		
		Der Text gliedert sich in drei Schichten.
		Zunächst werden die algebraischen Grundlagen der Hurwitz-Ordnung,
		endlicher Quotienten und Cayley-Graphen entwickelt.
		Dann werden die Operatoren als Faltungen auf endlichen Quotienten beschrieben
		und ihre spektralen Konsequenzen präzisiert.
		Schließlich wird erläutert, in welchem Sinn die normgesteuerten Operatoren
		als endliche Vorstufe heckeartiger Phänomene gelesen werden können,
		ohne dass dadurch bereits eine echte automorphe oder fargues--scholzesche
		Theorie behauptet würde.
		
		Der Artikel ist daher zugleich konstruktiv und disziplinierend:
		Er vertieft den mathematischen Kern der diskreten Konstruktion,
		markiert aber ebenso klar die Grenze zwischen gesicherter Spektraltheorie
		und weiterführender arithmetisch-geometrischer Perspektive.
	\end{abstract}
	
	\tableofcontents
	
	\section{Einleitung}
	
	Manchmal beginnt Mathematik nicht mit maximaler Allgemeinheit,
	sondern mit einer kontrollierten Verkleinerung.
	Gerade in der Zahlentheorie ist dies oft produktiv:
	Man ersetzt einen tiefen, schwer zugänglichen Zusammenhang
	durch ein endliches Modell, das strukturell ähnlich organisiert ist,
	aber vollständig explizit bleibt.
	
	Im vorliegenden Text geschieht dies ausgehend von der Hurwitz-Ordnung
	in der Hamiltonschen Quaternionenalgebra.
	Für jedes Niveau \(N\ge 1\) betrachten wir den endlichen Quotienten
	\[
	R_N := \OH / N\OH,
	\]
	sowie auf geeigneten Teilmengen definierte Faltungsoperatoren,
	deren Erzeugermengen über reduzierte Normbedingungen ausgewählt werden.
	Bereits auf dieser endlichen Ebene treten drei Motive zusammen:
	diskrete Dynamik, Spektraltheorie und Arithmetik.
	
	Der erste mathematische Gewinn ist elementar, aber präzise:
	Reguläre Graphen führen auf selbstadjungierte Operatoren,
	deren Spektrallücke exponentielle Durchmischung kontrolliert.
	Der zweite Gewinn ist konzeptionell:
	Sobald die Erzeugermengen aus Quaternionen fester reduzierter Norm
	stammen, entsteht eine Operatorfamilie, die formal an Hecke-Operatoren erinnert.
	Sie ist noch keine Hecke-Algebra im klassischen Sinn,
	aber sie besitzt dieselbe Organisationsidee:
	lokal standardisierte Schritte erzeugen global auswertbare Spektren.
	
	Gerade hier wird der Bezug zur modernen arithmetischen Geometrie interessant.
	Die Arbeiten von Peter Scholze und Laurent Fargues zeigen in großer Tiefe,
	dass arithmetische Information oft nicht als lose Sammlung lokaler Formeln,
	sondern in operatorischen, geometrischen und kategorialen Strukturen
	organisiert werden muss.
	Der vorliegende Text behauptet ausdrücklich \emph{nicht},
	dass das hier behandelte diskrete Quaternionenmodell bereits eine Instanz
	dieser Welt sei.
	Er beansprucht nur, ein endlicher Vorraum zu sein,
	in dem ein ähnliches Grundmuster sichtbar wird.
	
	\section{Die algebraische Ausgangslage}
	
	\subsection{Die Hamiltonsche Quaternionenalgebra}
	
	Wir arbeiten in der Quaternionenalgebra
	\[
	B := (-1,-1)_{\Q}
	= \Q \oplus \Q i \oplus \Q j \oplus \Q k,
	\]
	mit Relationen
	\[
	i^2=j^2=k^2=ijk=-1.
	\]
	Die Standardinvolution ist gegeben durch
	\[
	\overline{a+bi+cj+dk}=a-bi-cj-dk.
	\]
	Die reduzierte Spur und die reduzierte Norm lauten
	\[
	\trd(x)=x+\bar{x}, \qquad \Nrd(x)=x\bar{x}.
	\]
	Für \(x=a+bi+cj+dk\) ergibt sich explizit
	\[
	\Nrd(x)=a^2+b^2+c^2+d^2.
	\]
	
	\begin{bemerkung}
		Die reduzierte Norm ist multiplikativ:
		\[
		\Nrd(xy)=\Nrd(x)\Nrd(y).
		\]
		Damit besitzt \(B\) eine natürliche arithmetische Größenfunktion,
		die sowohl lokal als auch global zentral ist.
	\end{bemerkung}
	
	\subsection{Die Hurwitz-Ordnung}
	
	Wir betrachten die Hurwitz-Ordnung
	\[
	\OH := \Z\!\left[i,j,k,\frac{1+i+j+k}{2}\right].
	\]
	Sie ist eine maximale Ordnung in \(B\).
	Als \(\Z\)-Modul hat \(\OH\) Rang \(4\).
	
	\begin{lemma}
		Die Menge \(\OH\) ist eine \(\Z\)-Ordnung in \(B\), also ein
		endlich erzeugter \(\Z\)-Unterring, der \(B\) über \(\Q\) aufspannt.
	\end{lemma}
	
	\begin{proof}
		Dies ist Standard für die Hurwitz-Ordnung.
		Sie enthält \(1\), ist unter Addition und Multiplikation abgeschlossen
		und erzeugt \(B\) als \(\Q\)-Vektorraum.
	\end{proof}
	
	\begin{bemerkung}
		Für \(x\in \OH\) gilt \(\Nrd(x)\in \Z_{\ge 0}\).
		Diese ganzzahlige Norm ist der Ausgangspunkt für die spätere
		Konstruktion normgesteuerter Operatoren.
	\end{bemerkung}
	
	\subsection{Endliche Quotienten}
	
	Für \(N\ge 1\) definieren wir den Quotientenring
	\[
	R_N := \OH/N\OH.
	\]
	
	\begin{lemma}
		Als additive Gruppe gilt
		\[
		R_N \cong (\Z/N\Z)^4.
		\]
		Insbesondere hat \(R_N\) genau \(N^4\) Elemente.
	\end{lemma}
	
	\begin{proof}
		Da \(\OH\) als \(\Z\)-Modul frei vom Rang \(4\) ist,
		folgt
		\[
		\OH/N\OH \cong (\Z/N\Z)^4
		\]
		als additive Gruppen.
	\end{proof}
	
	\begin{bemerkung}
		Der Ring \(R_N\) ist endlich, also auch seine Einheitengruppe \(R_N^\times\).
		Damit steht eine endliche Bühne für diskrete Operatoren bereit.
	\end{bemerkung}
	
	\section{Cayley-Graphen und Faltungsoperatoren}
	
	\subsection{Von Erzeugermengen zu Graphen}
	
	Sei \(G_N\) eine endliche Gruppe, etwa \(R_N^\times\) oder
	eine geeignete durch Normbedingungen ausgewählte Untergruppe.
	Sei ferner \(\Sigma \subseteq G_N\) eine endliche Teilmenge mit
	\[
	\Sigma=\Sigma^{-1}, \qquad 1\notin \Sigma.
	\]
	
	\begin{definition}
		Der zu \((G_N,\Sigma)\) gehörige Cayley-Graph \(\mathrm{Cay}(G_N,\Sigma)\)
		hat Knotenmenge \(G_N\), und zwei Knoten \(g,h\in G_N\) werden verbunden,
		falls \(h=sg\) für ein \(s\in \Sigma\).
	\end{definition}
	
	Ist \(\Sigma\) symmetrisch, so handelt es sich um einen ungerichteten
	\(|\Sigma|\)-regulären Graphen.
	
	\subsection{Der Operator als Faltung}
	
	Der entscheidende Schritt besteht darin, den Adjazenzoperator nicht nur
	kombinatorisch, sondern algebraisch zu lesen.
	
	\begin{definition}
		Für eine endliche Gruppe \(G\) und eine symmetrische Teilmenge \(\Sigma\subseteq G\)
		definieren wir den Faltungsoperator
		\[
		T_\Sigma : \ell^2(G)\to \ell^2(G),\qquad
		(T_\Sigma f)(g)=\sum_{s\in \Sigma} f(s^{-1}g).
		\]
		Der normierte Markov-Operator ist
		\[
		P_\Sigma := \frac{1}{|\Sigma|}T_\Sigma.
		\]
	\end{definition}
	
	\begin{proposition}
		Ist \(\Sigma=\Sigma^{-1}\), so ist \(T_\Sigma\) selbstadjungiert
		bezüglich des Standardskalarprodukts auf \(\ell^2(G)\).
		Insbesondere ist das Spektrum von \(T_\Sigma\) reell.
	\end{proposition}
	
	\begin{proof}
		Für \(f,h\in \ell^2(G)\) gilt
		\[
		\langle T_\Sigma f,h\rangle
		=
		\sum_{g\in G}\sum_{s\in\Sigma}f(s^{-1}g)\overline{h(g)}.
		\]
		Mit der Substitution \(u=s^{-1}g\) folgt
		\[
		\langle T_\Sigma f,h\rangle
		=
		\sum_{u\in G}\sum_{s\in\Sigma}f(u)\overline{h(su)}.
		\]
		Da \(\Sigma=\Sigma^{-1}\), ist dies gleich
		\[
		\sum_{u\in G} f(u)\overline{\sum_{s\in\Sigma} h(s^{-1}u)}
		=
		\langle f,T_\Sigma h\rangle.
		\]
	\end{proof}
	
	\begin{bemerkung}
		Der übliche Adjazenzoperator eines Cayley-Graphen ist also ein
		spezifischer Faltungsoperator im Gruppenring \(\C[G]\).
		Diese Beobachtung ist für die spätere Hecke-Analogie zentral.
	\end{bemerkung}
	
	\section{Spektrum, Spektrallücke und Durchmischung}
	
	\subsection{Normierte Spektren}
	
	Sei \(G\) ein endlicher \(|\Sigma|\)-regulärer Graph mit Adjazenzoperator \(T_\Sigma\)
	und normiertem Operator \(P_\Sigma\).
	Dann gilt
	\[
	\|P_\Sigma\| \le 1.
	\]
	Schreiben wir das Spektrum von \(P_\Sigma\) als
	\[
	1=\mu_1 \ge \mu_2 \ge \cdots \ge \mu_n \ge -1,
	\]
	so ist die konstante Funktion stets Eigenfunktion zum Eigenwert \(1\).
	
	\begin{definition}
		Die Spektrallücke des normierten Operators \(P_\Sigma\) sei
		\[
		\Delta(\Sigma):=
		1-\max_{i\ge 2}|\mu_i|.
		\]
	\end{definition}
	
	\begin{satz}
		Ist \(\Delta(\Sigma)>0\), so gilt für jede Funktion \(f\in \ell^2(G)\)
		mit Mittelwert \(0\):
		\[
		\|P_\Sigma^m f\|_2 \le (1-\Delta(\Sigma))^m \|f\|_2
		\qquad (m\ge 0).
		\]
	\end{satz}
	
	\begin{proof}
		Auf dem orthogonalen Komplement der Konstanten ist die Operatornorm von
		\(P_\Sigma\) gleich dem größten Betrag eines nichttrivialen Eigenwertes.
		Daher
		\[
		\|P_\Sigma^m\|_{\ell^2_0\to \ell^2_0}
		=
		\left(\max_{i\ge 2}|\mu_i|\right)^m
		=
		(1-\Delta(\Sigma))^m.
		\]
	\end{proof}
	
	\begin{korollar}
		Eine positive Spektrallücke impliziert exponentielle Durchmischung des
		zugehörigen Random Walks und damit eine präzise Form diskreter
		algorithmischer Effizienz.
	\end{korollar}
	
	\begin{bemerkung}
		Mit \enquote{Effizienz} ist hier ausschließlich die Mischungsrate des
		diskreten dynamischen Systems gemeint.
		Physikalische oder thermodynamische Energiebegriffe sind damit
		nicht verbunden.
	\end{bemerkung}
	
	\section{Normgesteuerte Operatoren aus Quaternionen}
	
	\subsection{Norm-\texorpdfstring{$p$}{p}-Elemente}
	
	Bis hierhin war \(\Sigma\) eine beliebige symmetrische Erzeugermenge.
	Der arithmetische Gehalt beginnt erst, wenn \(\Sigma\) aus der Quaternionenalgebra
	selbst stammt.
	
	\begin{definition}
		Sei \(p\) eine Primzahl.
		Wir definieren die Menge der Norm-\(p\)-Elemente in \(\OH\) durch
		\[
		X_p := \{x\in \OH \;:\; \Nrd(x)=p\}.
		\]
	\end{definition}
	
	\begin{bemerkung}
		Da \(\Nrd\) multiplikativ ist, sind diese Elemente natürliche
		lokale Schrittgeneratoren eines festen arithmetischen Typs.
		Will man daraus Operatoren auf \(R_N\) oder \(R_N^\times\) bauen,
		muss man modulo \(N\OH\) reduzieren und zugleich Einheitenquotienten
		berücksichtigen, um Mehrfachzählungen zu kontrollieren.
	\end{bemerkung}
	
	\begin{definition}
		Eine Teilmenge \(\Sigma_{N,p}\subseteq R_N^\times\) heiße
		\emph{normkompatibel vom Niveau \(p\)}, wenn sie aus Reduktionsklassen
		von Elementen aus \(X_p\) erzeugt wird, modulo Einheiten und gegebenenfalls
		einer weiteren symmetrisierenden Äquivalenzrelation, so dass
		\[
		\Sigma_{N,p}^{-1}=\Sigma_{N,p}
		\]
		gilt.
	\end{definition}
	
	\begin{bemerkung}
		Hier liegt einer der mathematisch sensibelsten Punkte des gesamten Ansatzes.
		Für eine wirklich starke Theorie genügt nicht irgendeine normkompatible Wahl;
		notwendig wäre eine \emph{kanonische} Konstruktion der Mengen \(\Sigma_{N,p}\),
		die mit lokaler Algebra und Doppelnebenklassenformalismus verträglich ist.
	\end{bemerkung}
	
	\subsection{Die zugehörigen Operatoren}
	
	Zu einer normkompatiblen Menge \(\Sigma_{N,p}\) definieren wir
	\[
	T_{N,p}:=T_{\Sigma_{N,p}},\qquad
	P_{N,p}:=\frac{1}{|\Sigma_{N,p}|}T_{N,p}.
	\]
	
	\begin{proposition}
		Ist \(\Sigma_{N,p}\) symmetrisch, so ist \(T_{N,p}\) selbstadjungiert.
		Falls \(\Sigma_{N,p}\) die betrachtete Komponente erzeugt,
		ist der zugehörige Cayley-Graph regulär und zusammenhängend.
	\end{proposition}
	
	\begin{proof}
		Die Selbstadjungiertheit folgt wie zuvor aus der Symmetrie.
		Die Regulärität ist Cayley-graphisch unmittelbar.
		Zusammenhang folgt aus der Erzeugereigenschaft.
	\end{proof}
	
	\begin{bemerkung}
		Formal sammelt \(T_{N,p}\) also lokale arithmetische Schritte fester Norm
		zu einem globalen Operator. Genau hierin beginnt die Nähe zur Hecke-Philosophie.
	\end{bemerkung}
	
	\section{Lokale Quaternionenalgebra und Hecke-Hintergrund}
	
	\subsection{Lokale Struktur von \texorpdfstring{$B$}{B}}
	
	Die Hamiltonsche Quaternionenalgebra \(B=(-1,-1)_\Q\) ist global nicht trivial,
	aber lokal an fast allen Stellen gesplittet.
	
	\begin{satz}
		Für jede ungerade Primzahl \(p\) gilt
		\[
		B\otimes_\Q \Q_p \cong M_2(\Q_p).
		\]
		Dagegen ist \(B\) an \(p=2\) und an der reellen Stelle ramifiziert.
	\end{satz}
	
	\begin{bemerkung}
		Für \(p\neq 2\) ist also auch eine maximale Ordnung lokal isomorph zu
		\[
		\OH\otimes_\Z \Z_p \cong M_2(\Z_p).
		\]
		Dies ist der eigentliche Eintrittspunkt lokaler Hecke-Geometrie.
		Denn in der gesplitteten Situation erscheinen Normbedingungen als
		Determinantenbedingungen für Matrizen.
	\end{bemerkung}
	
	\subsection{Doppelnebenklassen als kontinuierliche Vorlage}
	
	Die klassische Hecke-Theorie für \(\PGL_2(\Q_p)\) arbeitet mit
	kompakt gestützten \(K\)-biinvarianten Funktionen,
	wobei
	\[
	K = \PGL_2(\Z_p)
	\]
	eine maximale kompakte Untergruppe ist.
	Der Hecke-Operator der Stufe \(p\) entspricht formal der Doppelnebenklasse
	\[
	K \diag(p,1) K.
	\]
	
	\begin{bemerkung}
		Die zentrale Idee lautet: Man mittelt nicht über beliebige Gruppenelemente,
		sondern über eine lokal kanonische Bahn fester Determinante.
		Gerade deshalb tragen Hecke-Operatoren tiefe arithmetische Information.
	\end{bemerkung}
	
	\begin{proposition}[heuristische Brücke]
		Sind die Mengen \(\Sigma_{N,p}\) tatsächlich aus Reduktionsklassen
		von Norm-\(p\)-Elementen konstruiert und lokal mit der gesplitteten Struktur
		\(\OH\otimes \Z_p \cong M_2(\Z_p)\) kompatibel, dann ist \(T_{N,p}\)
		als endliche Schattenversion eines Hecke-Typs zu lesen:
		nicht als echter Hecke-Operator, wohl aber als diskreter Quotientenoperator,
		der dieselbe lokale Organisationsidee nachahmt.
	\end{proposition}
	
	\begin{proof}[Kommentar statt Beweis]
		Hier liegt derzeit kein vollständiger mathematischer Beweis vor,
		sondern ein Programm.
		Die Aussage ist als präzise formulierte Heuristik zu verstehen,
		deren Belastbarkeit von einer kanonischen Konstruktion von \(\Sigma_{N,p}\),
		lokalen Identifikationen und Kompatibilitätsnachweisen abhängt.
	\end{proof}
	
	\section{Bruhat--Tits-Bäume und die geometrische Lesart}
	
	\subsection{Der lokale Baum}
	
	Für \(p\neq 2\) besitzt \(\PGL_2(\Q_p)\) seinen Bruhat--Tits-Baum \(\mathcal{T}_p\),
	einen \((p+1)\)-regulären Baum.
	Seine Knoten können als Homothetyklassen von \(\Z_p\)-Gittern in \(\Q_p^2\)
	aufgefasst werden.
	Nachbarschaft entspricht dann einem elementaren Index-\(p\)-Schritt.
	
	\begin{bemerkung}
		Die klassische Hecke-Nachbarschaftsrelation auf diesem Baum ist genau die
		lokale Geometrie, aus der viele arithmetische Graphkonstruktionen hervorgehen.
		Insbesondere entstehen Ramanujan-Graphen in berühmten Beispielen
		durch geeignete arithmetische Quotienten solcher Bäume.
	\end{bemerkung}
	
	\subsection{Diskrete Reduktion als endlicher Schatten}
	
	Der vorliegende Ansatz arbeitet nicht direkt mit \(\mathcal{T}_p\),
	sondern mit endlichen Quotienten \(R_N\).
	Gerade deshalb ist er schwächer, aber expliziter.
	
	\begin{hypothese}
		Es gibt eine präzise Konstruktion, in der die normkompatiblen Mengen
		\(\Sigma_{N,p}\) als endliche Reduktionsspuren lokaler Nachbarschaften
		des Bruhat--Tits-Baums verstanden werden können.
	\end{hypothese}
	
	\begin{bemerkung}
		Wäre eine solche Brücke hergestellt, dann erhielte der Operator \(T_{N,p}\)
		eine geometrische Lesart:
		Er wäre nicht mehr bloß Adjazenz auf einem künstlich gewählten Graphen,
		sondern der endliche Reflex einer lokal kanonischen \(p\)-Nachbarschaft.
	\end{bemerkung}
	
	\section{Spektrale Konsequenzen und mögliche Ramanujan-Perspektiven}
	
	\subsection{Vom endlichen Spektrum zur arithmetischen Schranke}
	
	Solange \(T_{N,p}\) nur ein selbstadjungierter Operator auf einem endlichen Graphen ist,
	kontrolliert seine Spektrallücke die Durchmischung.
	Entsteht er jedoch aus einer echten Hecke-Algebra,
	so tragen seine Eigenwerte potentiell automorphe Information.
	
	\begin{bemerkung}
		Die Ramanujan-Philosophie besagt grob,
		dass nichttriviale Eigenwerte solcher Operatoren
		durch das lokale Temperiertheitsverhalten
		automorpher Darstellungen kontrolliert werden.
		Im diskreten Modell wäre dies die Stelle,
		an der eine bloß graphentheoretische Spektrallücke
		in eine arithmetische Eigenwertschranke übergeht.
	\end{bemerkung}
	
	\begin{hypothese}
		Für geeignete Familien \((N,p)\) könnten die Operatoren \(T_{N,p}\)
		nach angemessener Normierung Spektren besitzen,
		deren nichttriviale Eigenwerte Schranken erfüllen,
		die an Ramanujan-artige Phänomene erinnern.
	\end{hypothese}
	
	\begin{bemerkung}
		Der vorliegende Text beweist eine solche Schranke nicht.
		Er markiert nur den mathematisch richtigen Ort,
		an dem eine solche Frage gestellt werden müsste.
	\end{bemerkung}
	
	\section{Langlands als Organisationsidee}
	
	Das Langlands-Programm ist in seiner vollen Gestalt zu groß,
	um hier dargestellt zu werden.
	Für den vorliegenden Zusammenhang genügt ein minimales,
	aber tragfähiges Strukturmotiv:
	lokale Daten werden durch Operatoren in globale Spektren übersetzt.
	
	Dieses Motiv erscheint im diskreten Modell in der Form
	\[
	\text{Norm-}p\text{-Elemente}
	\longrightarrow
	T_{N,p}
	\longrightarrow
	\Spec(T_{N,p})
	\longrightarrow
	\text{globale Mischungs- und Strukturinformation}.
	\]
	Natürlich ist dies nur ein endlicher Schatten der eigentlichen Theorie.
	An die Stelle automorpher Darstellungen treten Funktionen auf endlichen Gruppen;
	an die Stelle lokaler Korrespondenzen treten normgesteuerte Übergänge;
	an die Stelle geometrischer Kategorien treten Quotientengraphen.
	
	\begin{bemerkung}
		Gerade diese Schwäche ist aber erkenntnisfördernd.
		Denn sie zeigt, welche Teile der Organisationsidee bereits auf rein endlicher Ebene
		sichtbar werden und welche erst in echter geometrischer Tiefe entstehen.
	\end{bemerkung}
	
	\section{Der Scholze-Bezug in präziser mathematischer Form}
	
	Hier ist begriffliche Disziplin entscheidend.
	Peter Scholzes Arbeiten zu perfectoiden Räumen,
	prismatischen Methoden, geometrisierter lokaler Langlands-Korrespondenz,
	Wild-Betti-Garben und neuen Sechs-Funktoren-Formalismsen
	haben den begrifflichen Horizont der arithmetischen Geometrie tief verschoben.
	Dort wird arithmetische Information nicht mehr als lose Ansammlung
	punktweiser Invarianten behandelt,
	sondern in funktoriell organisierten geometrischen und kategorialen Räumen.
	
	Für das vorliegende Modell folgt daraus jedoch gerade \emph{nicht},
	dass es bereits ein Spezialfall dieser Theorien wäre.
	Die mathematisch korrekte Formulierung ist schärfer und bescheidener:
	
	\begin{quote}
		Das hier konstruierte Quaternionenmodell ist ein endlicher spektraler Proberaum.
		Es zeigt, wie lokal arithmetisch motivierte Schritte einen globalen Operator erzeugen,
		dessen Spektrum strukturelle Information trägt.
		In diesem methodischen, aber nicht technischen Sinn
		steht es unter einem Leitgedanken,
		der in den Arbeiten von Scholze und Fargues in weit tieferer Form verwirklicht wird.
	\end{quote}
	
	Es ist hilfreich, drei Ebenen streng auseinanderzuhalten:
	
	\begin{description}[leftmargin=2.3cm,style=nextline]
		\item[Stufe A.] Endliche Quotienten, Cayley-Graphen, Faltungsoperatoren,
		Spektrallücken, Random Walks.
		\item[Stufe B.] Hecke-artige Operatoren, lokale Quaternionenalgebren,
		Bruhat--Tits-Bäume, automorphe Spektren, Ramanujan-Phänomene.
		\item[Stufe C.] Fargues--Fontaine-Kurve, geometrisierte lokale Langlands-Korrespondenz,
		motivische Koeffizienten, Garbenkategorien, Sechs-Funktoren-Formalismen.
	\end{description}
	
	Der vorliegende Artikel gehört vollständig zu Stufe A,
	entwickelt eine präzisere Brücke zu Stufe B
	und nimmt Stufe C nur als fernes, aber instruktives Bezugssystem wahr.
	
	\begin{bemerkung}
		Unsauber wäre es, \(T_{N,p}\) bereits als Hecke-Operator
		im Sinn von Fargues--Scholze auszugeben oder
		den Quotientengraphen \(\mathrm{Cay}(G_N,\Sigma_{N,p})\)
		als diskrete Version der Fargues--Fontaine-Geometrie zu bezeichnen.
		Genau gegen eine solche Überdehnung richtet sich die saubere Architektur des Textes.
	\end{bemerkung}
	
	\section{Was formal gesichert ist -- und was nicht}
	
	\begin{proposition}
		Im Rahmen dieses Artikels sind die folgenden Aussagen formal gesichert:
		\begin{enumerate}[label=\arabic*.]
			\item Die Hurwitz-Ordnung \(\OH\) ist eine maximale \(\Z\)-Ordnung in \(B\).
			\item Für jedes \(N\ge 1\) ist \(R_N=\OH/N\OH\) ein endlicher Ring mit \(N^4\) Elementen.
			\item Jede symmetrische endliche Erzeugermenge \(\Sigma\) erzeugt einen endlichen
			regulären Cayley-Graphen und einen selbstadjungierten Faltungsoperator \(T_\Sigma\).
			\item Eine positive Spektrallücke impliziert exponentielle \(\ell^2\)-Durchmischung.
			\item Normgesteuerte Mengen \(\Sigma_{N,p}\) liefern eine arithmetisch motivierte
			Operatorfamilie \(T_{N,p}\), deren formale Struktur heckeartig gelesen werden kann.
		\end{enumerate}
	\end{proposition}
	
	\begin{proposition}
		Nicht bewiesen werden in diesem Text:
		\begin{enumerate}[label=\arabic*.]
			\item eine kanonische und vollständig funktorielle Konstruktion von \(\Sigma_{N,p}\),
			\item eine Identifikation von \(T_{N,p}\) mit klassischen Hecke-Operatoren,
			\item eine Ramanujan-Eigenschaft der entstehenden Graphen,
			\item ein Vergleich mit automorphen Spektren im strengen Sinn,
			\item eine Einbettung in perfectoide, prismatische oder garbentheoretische Geometrie,
			\item eine Realisierung der lokalen Langlands-Korrespondenz aus dem diskreten Modell.
		\end{enumerate}
	\end{proposition}
	
	\begin{bemerkung}
		Diese Negativliste ist kein Mangel,
		sondern eine methodische Stärke.
		Sie trennt bewiesene endliche Spektraltheorie
		von weiterführender arithmetischer Vision.
	\end{bemerkung}
	
	\section{Wie aus der Analogie mehr werden könnte}
	
	Soll aus der gegenwärtigen Heuristik eine belastbare arithmetische Theorie werden,
	dann wären mindestens die folgenden Schritte erforderlich:
	
	\begin{enumerate}[label=\arabic*.]
		\item Eine \emph{kanonische} Konstruktion der Mengen \(\Sigma_{N,p}\)
		aus Norm-\(p\)-Elementen der Hurwitz-Ordnung modulo Einheiten.
		\item Eine genaue Analyse der lokalen Ringe
		\[
		\OH\otimes \Z_\ell
		\]
		und ihrer Beziehung zu \(\PGL_2(\Q_\ell)\) sowie zu Bruhat--Tits-Bäumen.
		\item Der Nachweis, dass die Operatoren \(T_{N,p}\) in geeigneten Familien
		aus einer echten Hecke-Algebra stammen.
		\item Eine simultane spektrale Analyse dieser Operatorfamilien,
		idealerweise mit Vergleich zu automorphen oder Ramanujan-artigen Schranken.
		\item Eine Einbettung des diskreten Modells in eine moduli-,
		garben- oder korrespondenztheoretische Umgebung.
	\end{enumerate}
	
	\begin{bemerkung}
		Erst auf dieser Grundlage wäre ein stärkerer Langlands- oder Scholze-Bezug
		mathematisch voll belastbar.
		Vorher bleibt die Analogie produktiv, aber begrenzt.
	\end{bemerkung}
	
	\section{Schluss}
	
	Der eigentliche Gewinn des vorliegenden Ansatzes liegt nicht in einer
	spektakulären Großbehauptung,
	sondern in einer präzisen Zwischenebene:
	Er zeigt, dass man aus einer maximalen Quaternionenordnung
	durch endliche Reduktion, normgesteuerte Erzeugermengen und Faltungsoperatoren
	eine nichttriviale spektrale Welt erzeugen kann,
	in der lokale arithmetische Information global lesbar wird.
	
	Gerade dadurch gewinnt der Bezug zu moderner arithmetischer Geometrie an Schärfe.
	Nicht weil das Modell bereits geometric Langlands,
	Fargues--Fontaine oder Scholzes Sechs-Funktoren-Welt realisierte,
	sondern weil es in endlicher Form ein Motiv isoliert,
	das dort in hoher begrifflicher Dichte wiederkehrt:
	die Übersetzung lokaler Struktur in globale Operatorik.
	
	So gelesen ist das vorliegende Quaternionenmodell weder bloß
	elementare Graphentheorie noch schon tiefe Langlands-Geometrie.
	Es ist etwas Drittes:
	ein endlicher spektraler Vorraum,
	in dem die Frage nach echter arithmetischer Tiefe
	überhaupt erst in einer mathematisch kontrollierten Form gestellt werden kann.
	
	\begin{thebibliography}{99}
		
		\bibitem{ConwaySloane}
		J. H. Conway und N. J. A. Sloane,
		\emph{Sphere Packings, Lattices and Groups},
		3. Auflage, Springer, 1999.
		
		\bibitem{Lubotzky}
		A. Lubotzky,
		\emph{Discrete Groups, Expanding Graphs and Invariant Measures},
		Birkhäuser, 1994.
		
		\bibitem{SerreTrees}
		J.-P. Serre,
		\emph{Trees},
		Springer, 1980.
		
		\bibitem{Terras}
		A. Terras,
		\emph{Fourier Analysis on Finite Groups and Applications},
		Cambridge University Press, 1999.
		
		\bibitem{Vigneras}
		M.-F. Vignéras,
		\emph{Arithmétique des algèbres de quaternions},
		Springer, 1980.
		
		\bibitem{Titchmarsh}
		E. C. Titchmarsh,
		\emph{The Theory of the Riemann Zeta-Function},
		2. Auflage, herausgegeben von D. R. Heath-Brown,
		Oxford University Press, 1986.
		
		\bibitem{ScholzeMotivic}
		P. Scholze,
		\emph{Geometrization of the local Langlands correspondence, motivically},
		arXiv:2501.07944, 2025.
		
		\bibitem{ScholzeWildBetti}
		P. Scholze,
		\emph{Wild Betti sheaves},
		arXiv:2505.24599, 2025.
		
		\bibitem{ScholzeSixFunctor}
		P. Scholze,
		\emph{Six-Functor Formalisms},
		arXiv:2510.26269, 2025.
		
		\bibitem{AnschutzBoscoLeBrasRodriguezCamargoScholze}
		J. Anschütz, G. Bosco, A.-C. Le Bras, J. E. Rodríguez Camargo und P. Scholze,
		\emph{Analytic de Rham stacks of Fargues--Fontaine curves},
		arXiv:2510.15196, 2025.
		
		\bibitem{FarguesScholze}
		L. Fargues und P. Scholze,
		\emph{Geometrization of the local Langlands correspondence},
		Buchmanuskript / Preprint, 2021 ff.
		
	\end{thebibliography}
	
\end{document}